% Generated draft: requires mathematical and source review.
\documentclass[11pt]{article}
\usepackage[margin=25mm]{geometry}
\usepackage{fontspec}
\setmainfont{Latin Modern Roman}
\usepackage{amsmath,amssymb,mathtools,unicode-math}
\setmathfont{Latin Modern Math}
\usepackage{enumitem,needspace,xurl,hyperref}
\setlength{\parindent}{0pt}
\setlength{\parskip}{5pt}
\setlength{\emergencystretch}{3em}
\newcommand{\N}{\mathbb{N}}
\newcommand{\Z}{\mathbb{Z}}
\newcommand{\Q}{\mathbb{Q}}
\newcommand{\R}{\mathbb{R}}
\newcommand{\C}{\mathbb{C}}
\newcommand{\F}{\mathbb{F}}
\newcommand{\A}{\mathbb{A}}
\newcommand{\PP}{\mathbb{P}}
\newcommand{\E}{\mathbb{E}}
\newcommand{\eps}{\varepsilon}
\newcommand{\dd}{\,\mathrm{d}}
\begin{document}
\section*{UnsolvedMath ID: 30006625}
\textbf{Source problem code:} catalogue entry
\subsection{Definitions and mathematical statement}
For metric spaces \((Y,d_Y)\) and \((X,d_X)\), a quasi-isometry is a map \(f:Y\to X\) for which there exist \(L\geq 1\) and \(A\geq 0\) such that
\[
L^{-1}d_Y(y,z)-A\leq d_X(f(y),f(z))\leq Ld_Y(y,z)+A
\quad(y,z\in Y),
\]
and every point of \(X\) lies at distance at most \(A\) from \(f(Y)\).

For a chosen metric realization \((G,d)\) of a group, let \(\omega\) be a nonprincipal ultrafilter on \(\N\), let \(s_n>0\) satisfy \(s_n\to\infty\), and let \(b_n\in G\) be basepoints. Define
\[
\operatorname{Cone}_{\omega}(G,d;(b_n),(s_n))
=
\left\{(g_n):\sup_n\frac{d(g_n,b_n)}{s_n}<\infty\right\}\big/\sim,
\]
where
\[
(g_n)\sim(h_n)
\iff
\lim_{n\to\omega}\frac{d(g_n,h_n)}{s_n}=0.
\]
The cone metric is
\[
d_{\omega}([(g_n)],[(h_n)])
=\lim_{n\to\omega}\frac{d(g_n,h_n)}{s_n}.
\]
Two metric spaces are bi-Lipschitz equivalent if there is a bijection \(F\) between them and a constant \(B\geq1\) such that
\[
B^{-1}d(u,v)\leq d(F(u),F(v))\leq B d(u,v)
\quad\text{for all }u,v.
\]

\textbf{Source question.} Within the class of coarse median groups, characterize those groups that are quasi-isometric to a hierarchically hyperbolic space, using their asymptotic cones. Thus the requested characterization should give a condition on asymptotic cones equivalent to
\[
\exists X\;\bigl(X\text{ is a hierarchically hyperbolic space and }G\text{ is quasi-isometric to }X\bigr).
\]
The terms ``coarse median group'' and ``hierarchically hyperbolic space'' are retained as the technical notions used by the supplied source; their defining axioms are not supplied in the record. The record also does not specify the metric convention for groups, finite generation, or any rank restriction. The cone construction above applies after a metric has been chosen; it does not resolve those missing hypotheses.

\textbf{Additional formulation in the curated statement.} The supplied clean statement asks whether the preceding existence condition is equivalent to
\[
\begin{gathered}
\text{for every nonprincipal }\omega,\ \text{every }(s_n)\text{ with }s_n>0,\ s_n\to\infty,\ \text{and every }(b_n)\subseteq G,\\
\exists U\;\bigl(U\text{ is a universal real cubing and }
\operatorname{Cone}_{\omega}(G,d;(b_n),(s_n))
\text{ is bi-Lipschitz equivalent to }U\bigr).
\end{gathered}
\]
Here the space \(U\) and the bi-Lipschitz constant may depend on the cone; no uniform constant is specified. This display makes explicit a possible reading of ``each nonprincipal asymptotic cone.'' However, the record does not define ``universal real cubing,'' its universality category, or any rank or cardinality parameters. Moreover, this specific biconditional does not occur in the supplied original excerpt. Consequently it is an unresolved curated formulation, not a fully specified equivalence attributable to the excerpt without further source review.
\subsection{Short English statement}
The source asks for an asymptotic-cone characterization of coarse median groups that are quasi-isometric to hierarchically hyperbolic spaces. The curated statement additionally proposes that every cone be bi-Lipschitz equivalent to a universal real cubing. That proposal requires clarification of the group metric, the meaning of universality, and its support in the original report.
\subsection{Sources}
\begin{itemize}
\item Local UnsolvedMath registry, permanent ID \(30006625\), source code \texttt{catalogue entry}, title: \emph{Asymptotic-Cone Characterization of Hierarchically Hyperbolic Spaces}. The supplied record belongs to the Oberwolfach Reports Open Problems collection and gives proposed year 2026; no proposer is identified.
\item Supplied report citation: \emph{Median Geometry and Applications} (2026), Oberwolfach Reports, pp. 487--540, DOI \texttt{10.4171/owr/2026/8}. Supplied source URL: \url{https://doi.org/10.4171/owr/2026/8}. The supplied literature assessment identifies the question as Problem 13(ii).
\item Supplied literature source, labelled ``OWR report'': \url{https://ems.press/content/serial-article-files/53603}.
\item Literature checked at: \texttt{2026-08-22}. This date and assessment are supplied dated curation, not independently verified. The curation labels the problem open, describes the asymptotic-cone characterization as a question in the report, and reports that no resolution was found. Its background expressly requests independent verification of time-sensitive status and cited claims.
\item The registry labels the statement \texttt{corrected\_verified} and says that only formula delimiters, notation, grammar, and typography were repaired while preserving mathematical content and scope. These are supplied editorial claims; the discrepancy between the original excerpt and the more specific curated statement remains unresolved.
\end{itemize}
\end{document}
